%% ma: L10-B02-0040
%% lop: 10
%% chuong: 1
%% bai: 2
%% dang: 10.2.9
%% loai: TLN
%% muc: TH
%% dap_an: "5"
%% nguon: "Ảnh giáo viên gửi 10/10/2026 (ThS. Nguyễn Hoàng Việt, Chương 2 Tập hợp), Đề 2 (đề thứ hai, không ghi tiêu đề), Phần 3 Câu 18"
%% kien_thuc: "Bài 2 (tập hợp và các phép toán trên tập hợp)"
%% trang_thai: chinh-thuc
\begin{ex}
Lớp $10$A có $45$ học sinh trong đó có $25$ em học giỏi môn Toán, $23$ em học giỏi môn Lý, $20$ em học giỏi môn Hóa, $11$ em học giỏi cả môn Toán và môn Lý, $8$ em học giỏi cả môn Lý và môn Hóa, $9$ em học giỏi cả môn Toán và môn Hóa. Hỏi lớp $10$A có bao nhiêu bạn học giỏi cả ba môn Toán, Lý, Hóa? (biết rằng mỗi học sinh trong lớp học giỏi ít nhất một trong ba môn Toán, Lý, Hóa).
\shortans{5}
\loigiai{Theo công thức $n(A\cup B\cup C)=n(A)+n(B)+n(C)-n(A\cap B)-n(B\cap C)-n(C\cap A)+n(A\cap B\cap C)$: $45=25+23+20-11-8-9+x=40+x$, suy ra $x=5$.}
\end{ex}
